% =====================================================================
%  RQ3 v2 — PAPER-READY CONTENT (real numbers, DiffuCoder-7B-Base)
%  Universe: 230 RefineID snippets / 1055 [MASK] sites / 196 names / 64 packages.
%  Figures already rendered to figures/new/rq3/:
%     fig_rq3_exp1_depth_probe.pdf   fig_rq3_exp2_dcl.pdf   fig_rq3_regime_inversion.pdf
%  Copy each block into the corresponding place in the thesis. All numbers below
%  are taken verbatim from results/rq3_probe/exp{1,2}_curves.json.
% =====================================================================


% ---------------------------------------------------------------------
% (A) ABSTRACT — replace the current RQ3 sentence with this
% ---------------------------------------------------------------------
% OLD: "Probing the internal states of DiffuCoder-7B shows why: the signal that
% tells a bad name from a good one appears only in the last few layers and the
% last few denoising steps, after the unmasking schedule has already confirmed
% most of its predictions."
%
% NEW:
A linear probe on DiffuCoder-7B-Base shows why. Whether an identifier fits its
surrounding code is linearly decodable from the residual stream, but the
genuinely contextual component of that signal (intact minus scrambled-context
decodability) is near zero through the first two-thirds of the network and rises
only in the final layers, reaching an AUC of $0.26$ at the output. During
generation the model carries no early readable signal of whether its own fill
will be correct: commitment concentrates in the final third of the denoising
schedule, while a probe for the model's eventual success stays near chance
(peak ROC $0.63$) up to that point. The discriminative information therefore
exists but emerges too late, and is not available as an actionable signal before
the schedule commits, which is why providing the rename positions as masks
bypasses the problem and dLLMs work as filling engines but not standalone agents.


% ---------------------------------------------------------------------
% (B) THE RQ3 SECTION  (replaces the old four-probe Section VI)
% ---------------------------------------------------------------------
\section{RQ3: What the Model Represents About Name Fit}
\label{sec:rq3}
RQ1 and RQ2 characterized DiffuCoder from the outside, by what it produces at
masked positions. A model that recovers the developer's exact identifier a third
of the time must compute something internally that separates the names it keeps
from the names it would replace. We probe that internal state on $n=230$ RefineID
snippets ($1055$ sites, $196$ distinct names, $64$ source packages), using
DiffuCoder-7B-Base because its inference path exposes per-layer hidden states and
the unmasking schedule through a stable interface.

\subsection{Name fit is linearly decodable, contextually, and only in deep layers}
\label{sec:rq3-depth}
For each snippet we clamp the target identifier to the developer's name (the
\emph{good} condition) or to a length- and sub-word-count--matched
\emph{misplaced} real identifier drawn from a different snippet (the \emph{bad}
condition), holding the surrounding code byte-identical between the two. This
replaces the trivial ``real word versus short throw-away token'' contrast---which
a probe solves at $\mathrm{AUC}=1.0$ by reading token length alone
(Appendix~\ref{app:rq3})---with ``the right name versus an equally-shaped wrong
name in this context.'' We read the residual stream at the target position across
all $29$ layers and train an $\ell_2$ logistic probe at each layer, evaluated
under leave-one-package-out cross-validation with snippet-cluster bootstrap
confidence intervals. The headline is the \emph{contextual} decodability,
$\mathrm{AUC}_{\text{ctx}}(\ell)=\mathrm{AUC}_{\text{intact}}(\ell)-
\mathrm{AUC}_{\text{scrambled}}(\ell)$, where the scrambled control clamps the
same tokens but shuffles the surrounding code; any token-intrinsic signal
survives scrambling and cancels, leaving only what depends on the intact context.

Raw decodability is high but largely token-intrinsic: with intact context the
probe separates good from misplaced names at $\mathrm{AUC}=0.956$
($95\%\,\mathrm{CI}\,[0.939,0.972]$) at the output layer, yet the scrambled
control still reaches $0.693$, because the developer's name shares lexical
material with its own snippet. The contextual difference is therefore $0.263$,
and---unlike the raw curve, which plateaus near $0.99$ by the middle of the
stack---it is flat at $\approx0.10$ through the first two-thirds of the network
(early-third mean $0.103$) and climbs monotonically over the final third
(late-third mean $0.186$) to its peak of $0.263$ at the output
(Fig.~\ref{fig:rq3-exp1}a). The scrambled curve \emph{falls} from $0.87$
mid-stack to $0.69$ at the output over the same range, i.e.\ deep layers
increasingly depend on the actual surrounding structure rather than the token
alone. The signal is not an artifact: the sub-word-count baseline is exactly
$0.500$, the random-label control task is $0.534$ (selectivity $0.422$), the
layer-0 embedding is $0.301$, and grouping by snippet rather than by package
inflates the AUC by only $0.006$, bounding codebase leakage
(Table~\ref{tab:rq3-exp1}). The probe's decision axis cleanly separates the two
classes at the output layer (Fig.~\ref{fig:rq3-exp1}b).

\paragraph{Cross-model replication}
The same depth profile holds on DreamCoder-7B, a second masked-diffusion code
model: the contextual signal is flat at $\approx0.11$ through the first two-thirds
and rises to $0.297$ at the output, tracking the DiffuCoder-7B-Base curve almost
exactly (Fig.~\ref{fig:rq3-compare}), with every control passing (intact $0.978$,
scrambled $0.681$, sub-word-count baseline $0.500$, selectivity $0.511$, leakage
gap $0.006$). The late-emerging, deep-layer name-fit representation is therefore a
property of masked discrete diffusion rather than a single checkpoint.

\subsection{During generation, success is not decodable before commitment}
\label{sec:rq3-time}
The depth analysis shows where the fit signal lives; the schedule analysis asks
whether it arrives in time to act. We run normal generation (target masked,
filled over the $32$-step denoising schedule, EM rate $0.367$) and, at each step,
train a probe to predict whether the model's \emph{own} final fill will be exact,
evaluated only on positions still masked at that step so the already-committed
token cannot leak into the representation. We overlay this detection curve with
the commitment CDF (first step at confidence $\geq0.8$).

Commitment concentrates in the final third of the schedule: $25\%$ of positions
are committed by step $21$, $50\%$ by step $26$, and $75\%$ by step $29$. The
detection probe, however, stays near chance throughout: ROC rises from
$\approx0.50$ only to a peak of $0.63$ at step $22$ and never reaches $0.70$,
and PR-AUC sits at the $0.367$ base rate (Fig.~\ref{fig:rq3-exp2}). The result is
robust to the strictness of the success label: grading success by the CoReFusion
Identifier Score---which credits canvas-duplicated fills such as
\texttt{request}$\rightarrow$\texttt{requestRequest}---lifts the detection peak
only from $0.63$ to $0.67$ and leaves it at the same step ($22$, with $31\%$ of
positions already locked), so the absence of an early success signal is not an
artifact of exact match. The model therefore commits fills without an early,
internally-readable signal of whether they are correct.

\subsection{Why rare names are hardest}
\label{sec:rq3-regime}
Stratifying target positions by the rank of the ground-truth name in the output
distribution explains where the late signal fails to matter. In the
\textsc{HighConfident} regime the developer's name sits at median rank $77$ and an
injected smell at $492$; in \textsc{Uncertain} the gap narrows ($524$ vs $676$);
in \textsc{RareConfident} the ordering \emph{inverts}---the generic smell sits at
median rank $733$ while the developer's project-specific name sits at $10{,}989$
(Fig.~\ref{fig:rq3-regime}). These rare, project-specific identifiers are exactly
the ones a late-arriving, weak fit signal cannot rescue, linking RQ3 back to the
RQ2 failure mode. (The context-sensitivity analysis $\Delta H$ is in
Appendix~\ref{app:rq3}.)


% ---------------------------------------------------------------------
% (C) FIGURES
% ---------------------------------------------------------------------
\begin{figure}[t]
  \centering
  \includegraphics[width=\linewidth]{figures/new/rq3/fig_rq3_exp1_depth_probe.pdf}
  \caption{RQ3 Exp1. (a) Contextual name-fit decodability
  $\mathrm{AUC}_{\text{intact}}-\mathrm{AUC}_{\text{scrambled}}$ versus
  transformer layer (leave-one-package-out CV, snippet-cluster bootstrap band),
  with the intact, scrambled, length, and random-label-control references; the
  contextual signal is flat through the first two-thirds and rises in the final
  third (shaded) to $0.26$. (b) Out-of-fold projection of the layer-28 probe,
  developer name versus length-matched misplaced name (AUC $0.96$); this replaces
  the unsupervised UMAP view with the probe's own decision axis. $n=230$ snippets
  / $1055$ sites / $64$ packages, DiffuCoder-7B-Base.}
  \label{fig:rq3-exp1}
\end{figure}

\begin{figure}[t]
  \centering
  \includegraphics[width=\linewidth]{figures/new/rq3/fig_rq3_exp2_dcl.pdf}
  \caption{RQ3 Exp2. Per-step detection ROC and PR-AUC for ``will the model's own
  fill be exact,'' evaluated on still-masked positions, overlaid with the
  commitment CDF. Commitment completes in the final third (median step $26$),
  while detection never clears ROC $0.70$ (peak $0.63$ at step $22$) and PR-AUC
  stays at the $0.37$ base rate: the success signal does not precede commitment.}
  \label{fig:rq3-exp2}
\end{figure}

\begin{figure}[t]
  \centering
  \includegraphics[width=0.72\linewidth]{figures/new/rq3/fig_rq3_regime_inversion.pdf}
  \caption{RQ3 output-distribution view. Median rank of the developer's name
  versus an injected smell across confidence regimes; the ordering inverts in
  \textsc{RareConfident}, where the model ranks a generic smell above the
  developer's rare project-specific name.}
  \label{fig:rq3-regime}
\end{figure}

\begin{figure}[t]
  \centering
  \includegraphics[width=0.8\linewidth]{figures/new/rq3/fig_rq3_compare_models.pdf}
  \caption{RQ3 cross-model replication. Contextual name-fit decodability
  ($\mathrm{AUC}_{\text{intact}}-\mathrm{AUC}_{\text{scrambled}}$) versus
  normalized depth for DiffuCoder-7B-Base and DreamCoder-7B; both stay near zero
  through the first two-thirds and rise in the final third (shaded), so the
  deep-layer name-fit signal is shared across the masked-diffusion family.}
  \label{fig:rq3-compare}
\end{figure}


% ---------------------------------------------------------------------
% (D) RESULTS TABLE — Exp1 decodability and controls
% ---------------------------------------------------------------------
\begin{table}[t]
  \centering
  \caption{RQ3 Exp1: name-fit decodability and controls at the output layer
  (DiffuCoder-7B-Base, leave-one-package-out CV, $230$ snippets / $64$ packages).}
  \label{tab:rq3-exp1}
  \begin{tabular}{lc}
    \toprule
    Quantity & ROC-AUC \\
    \midrule
    Intact context (good vs.\ misplaced) & $0.956\;[0.939,0.972]$ \\
    Scrambled context (token-intrinsic floor) & $0.693$ \\
    \textbf{Contextual (intact $-$ scrambled)} & $\mathbf{0.263}$ \\
    \midrule
    Sub-word-count-only baseline & $0.500$ \\
    Random-label control task & $0.534$ \\
    Layer-0 embedding baseline & $0.301$ \\
    Leakage gap (snippet $-$ package CV) & $0.006$ \\
    Throw-away smell vocab (length-shortcut ceiling) & $1.000$ \\
    \bottomrule
  \end{tabular}
\end{table}


% ---------------------------------------------------------------------
% (E) METHODOLOGY ADDITION  (put in the Methodology chapter, RQ3 protocol)
% ---------------------------------------------------------------------
% \subsubsection{RQ3 internal-state protocol}
We probe DiffuCoder-7B-Base on the $230$-snippet, $1024$-token-filtered RefineID
split. \emph{Depth probe.} For each snippet we form a good/bad pair that differs
only at the target identifier: the good condition clamps the developer's name,
the bad condition a misplaced real identifier from another snippet matched on
sub-word count (char-length bucket as fallback), so length, frequency and
foreignness are controlled by construction. We read the residual stream at the
target's first sub-token for layers $0$--$28$ and fit an $\ell_2$ logistic probe
per layer. We report $\mathrm{AUC}_{\text{intact}}-\mathrm{AUC}_{\text{scrambled}}$
(scrambled = same tokens, permuted context) so token-intrinsic separability
cancels, against a sub-word-count baseline, a layer-0 embedding baseline, and a
Hewitt--Liang random-label control task. Generalization uses leave-one-package-out
GroupKFold over the $64$ packages (the gap to snippet-level grouping is reported
as the leakage estimate), and all confidence intervals are snippet-cluster
bootstraps. \emph{Schedule probe.} In normal generation we capture, per denoising
step, the hidden state at each target position and whether it is still masked, and
train a per-step probe (on still-masked positions only) for whether the model's
own final fill is exact, overlaid with the commitment CDF (first step at
confidence $\geq0.8$).


% ---------------------------------------------------------------------
% (F) CONCLUSION — replace the RQ3 sentence in Section IX
% ---------------------------------------------------------------------
Inside DiffuCoder-7B-Base, a linear probe shows that whether an identifier fits
its context is decodable from the residual stream---the genuinely contextual
component (intact minus scrambled-context) reaches AUC $0.26$ and emerges only in
the final third of the layers, above length, embedding, and control-task
baselines and robust to held-out packages---yet during generation the model
carries no early readable signal of its own success: commitment concentrates in
the final third of the denoising schedule while a success probe stays near chance
(peak ROC $0.63$) until then (RQ3). The discriminative information exists but
arrives too late to act on, which is why dLLMs are effective multi-site filling
engines when rename positions are given but not standalone refactoring agents.


% ---------------------------------------------------------------------
% (G) THREATS TO VALIDITY  (additions / replacements for the RQ3 items)
% ---------------------------------------------------------------------
% Internal: The RQ3 bad class is necessarily synthetic, since RefineID contains
% no natural pre-rename names; we mitigate the obvious-token confound by matching
% the bad name on sub-word count and by reporting the contextual double difference
% rather than raw AUC. The fixed $k{=}2$ generation canvas can duplicate short
% identifiers (e.g.\ "request" -> "requestRequest"), which depresses the Exp2 EM
% label; we verified the near-chance detection result is robust to a graded
% CIS-based success label (detection peak 0.63 -> 0.67, same late step 22).
% External: the depth probe replicates on DreamCoder-7B (Fig. ref{fig:rq3-compare}),
% so the deep-layer name-fit signal is not specific to DiffuCoder-7B-Base; the
% internal-state findings still use two masked-diffusion code models on 230
% snippets (64 packages), and other languages or generation configs may differ.
% Conclusion: per-layer and per-step AUCs are many correlated tests; we report the
% full curve with simultaneous bootstrap bands and a pre-registered late-vs-early-
% third contrast rather than a single selected cell, and grouping is at the
% package level with snippet-cluster bootstraps to avoid optimistic intervals.


% ---------------------------------------------------------------------
% (H) APPENDIX pointer
% ---------------------------------------------------------------------
% \section{RQ3 supplementary}\label{app:rq3}
% - Per-layer intact / scrambled / contextual curves (full table from exp1_curves.json).
% - Throw-away smell-vocab probe (AUC 1.000) demonstrating the length shortcut the
%   matched-misplaced design removes.
% - Context-sensitivity Delta H by regime (HighConfident 2.01, RareConfident 2.23,
%   Uncertain 2.42; 0.41-bit spread).
% - 3x3 t_detect x t_commit sensitivity surface (all censored: detection never
%   clears the thresholds).
